\section{A common realization and the physical renewal identity}
\label{sec:common-renewal}
The preceding moment estimate gives a common realization of the
unbounded record twists on a scale of Lebesgue spaces. It is useful to
construct this realization before seeking an anisotropic spectral
estimate: the collision and induced operators, their damping variables,
and their chronological composition can then be identified without
regularity assumptions on an indicator multiplier. The renewal algebra
is classical; compare \cite{GouezelRenewal}. The assertion here is its
exact realization for the moving physical first-return record.

\subsection{First-return operators on the common collision space}
Use $M=\T\times(-1,1)$ and $\dd\nu=(4\pi)^{-1}\dd\alpha\dd p$.
Let $P_R$ be multiplication by $\one_{Y_R^*}$ and let $E_R=I-P_R$.
Recall locally that
\[
 c_* = \nu(Y_R^*)=\frac{91}{10000\pi},\qquad
 \dd\nu_R^*=c_*^{-1}\one_{Y_R^*}\dd\nu.
\]
On $L^1(M,\nu)$, define
\[
 (\mathcal A_{R,\omega}h)(y)
   =e^{i\omega\cdot f_R(T_R^{-1}y)}h(T_R^{-1}y),\qquad
 f_R=(\kappa_{{\rm coll},1},\kappa_{{\rm coll},2},1,\tau_R),
\]
for real $\omega=(u_1,u_2,s,b)$. Collision singularities are discarded
as null sets. These are isometries. The operators $P_R,E_R$ are
contractions on every $L^p$, although no such assertion about their
action on anisotropic distribution spaces is needed or made here.

\begin{theorem}[Exact damped renewal realization]
\label{thm:common-renewal}
For $m\ge1$ set
\begin{equation}\label{eq:first-return-operator}
 \mathcal R_{m,R,\omega}
 =P_R\mathcal A_{R,\omega}
       (E_R\mathcal A_{R,\omega})^{m-1}P_R,
 \qquad
 \mathcal R_R(z,\omega)=\sum_{m\ge1}z^m\mathcal R_{m,R,\omega}.
\end{equation}
For $|z|<1$ the series converges in operator norm on $L^1(M,\nu)$ and
\begin{equation}\label{eq:induced-damped-bound}
 \|\mathcal R_R(z,\omega)h\|_1\le |z|\|P_Rh\|_1.
\end{equation}
On the closed subspace $P_RL^1$, with identity $P_R$, one has
\begin{equation}\label{eq:physical-renewal-identity}
 P_R(I-z\mathcal A_{R,\omega})^{-1}P_R
       =(I-\mathcal R_R(z,\omega))^{-1}.
\end{equation}
Moreover,
\begin{equation}\label{eq:schur-return-operator}
 \begin{split}
 \mathcal R_R(z,\omega)={}&zP_R\mathcal A P_R\\
 &+z^2P_R\mathcal A E_R
 (I-zE_R\mathcal A E_R)^{-1}
             E_R\mathcal A P_R,
 \end{split}
\end{equation}
where $\mathcal A=\mathcal A_{R,\omega}$ and the middle inverse acts
on $E_R L^1$.

For $|z|\le1$, the strongly defined boundary series represents the
actual first-return twist: for $y\in Y_R^*$ and $x=(F_R^*)^{-1}y$,
\begin{equation}\label{eq:actual-induced-realization}
 (\mathcal R_R(z,\omega)h)(y)
       =z^{r_R^*(x)}e^{i\omega\cdot G_R(x)}h(x).
\end{equation}
It is zero off $Y_R^*$. In particular, for every $n\ge1$,
\begin{equation}\label{eq:actual-record-pairing}
 c_*^{-1}\int_M\mathcal R_R(z,\omega)^n\one_{Y_R^*}\dd\nu
       =\int_{Y_R^*}z^{N_{n,R}}e^{i\omega\cdot J_{n,R}}\dd\nu_R^*.
\end{equation}
\end{theorem}
\begin{proof}
The collision map is invertible and preserves $\nu$. Poincar\'e
recurrence in both time directions makes $F_R^*$ invertible and
measure-preserving on $Y_R^*$, modulo null sets. In the product
\eqref{eq:first-return-operator}, the initial and final states lie in
$Y_R^*$ and the $m-1$ intervening states lie outside it. Thus its initial
domain is exactly $\{r_R^*=m\}$. Its weight is the exponential of the
sum of the $m$ actual collision records, hence $e^{i\omega\cdot G_R}$.
These initial domains partition $Y_R^*$ up to a null set. Their images
under the corresponding return branches also partition $Y_R^*$, by
invertibility of the induced map. Consequently
\[
 \sum_{m\ge1}|z|^m\|\mathcal R_{m,R,\omega}h\|_1
       =\int_{Y_R^*}|z|^{r_R^*}|h|\dd\nu
       \le |z|\|P_Rh\|_1 \quad (|z|\le1).
\]
Each summand has operator norm at most one, so for $|z|<1$ the
operator-norm series converges geometrically. At the boundary the last
display proves strong convergence for each $h$ and gives
\eqref{eq:actual-induced-realization}. It does not assert operator-norm
convergence of the boundary series.

Let $\mathcal U_0=P_R$ and $\mathcal U_m=P_R\mathcal A^mP_R$ for
$m\ge1$. Partition a trajectory with endpoints in $Y_R^*$ according to
its first return time. Chronological composition, with the first segment
on the right, gives
\[
 \mathcal U_m=\sum_{r=1}^m\mathcal U_{m-r}\mathcal R_{r,R,\omega}.
\]
All generating series converge in operator norm for $|z|<1$. Summing
the identity gives $\mathcal U(z)=P_R+\mathcal U(z)\mathcal R_R(z,\omega)$.
The inverse of $I-\mathcal R_R$ exists by
\eqref{eq:induced-damped-bound}, and
$\mathcal U(z)=P_R(I-z\mathcal A)^{-1}P_R$, proving
\eqref{eq:physical-renewal-identity}. Separating $m=1$ from the series
and summing the intervening outside-section iterates proves
\eqref{eq:schur-return-operator}. All these inverses are used only in
$|z|<1$.

Finally iterate \eqref{eq:actual-induced-realization} and change
variables by $(F_R^*)^n$. Its collision counts and record sums are
$N_{n,R}$ and $J_{n,R}$, respectively. Division by $c_*$ yields
\eqref{eq:actual-record-pairing}. No independence of successive
excursions enters the identity.
\end{proof}

\subsection{A fixed complex neighborhood for the unbounded twists}
To cross the damping boundary at $z=1$, use $z=e^w$. Define the operator
directly by the actual $n$th return, rather than by an unproved
continuation of a resolvent. For $y\in Y_R^*$ let $x=(F_R^*)^{-n}y$ and set
\begin{equation}\label{eq:block-scale-operator}
 (\mathcal T_{n,R}(w,\omega)h)(y)
       =e^{wN_{n,R}(x)+i\omega\cdot J_{n,R}(x)}h(x),
\end{equation}
with zero extension to $M\setminus Y_R^*$. Here $w\in\mathbb C$ and
$\omega\in\mathbb C^4$. Holomorphy is stated on this universal cover.
The formulas are invariant under real shifts by $2\pi\mathbb Z^3$
in the first three frequency coordinates. Local complex charts about
real torus frequencies are restrictions of the same formulas; no
additional global complex-torus construction is being asserted. The
symbols $J_{n,R}$ and $G_R$ denote four real coordinates, with the
displacement contributing two coordinates.

\begin{theorem}[Holomorphic realization on a common Lebesgue scale]
\label{thm:common-scale-holomorphy}
Fix $1\le q<p\le\infty$ and put $t=(q^{-1}-p^{-1})^{-1}$. Let
\[
 \sigma(w,\omega)=|\operatorname{Re}w|+D|\operatorname{Im}\omega|_\infty,
 \qquad \Omega_{p,q}=\{(w,\omega):t\sigma(w,\omega)<c\}.
\]
For every $n\ge1$, \eqref{eq:block-scale-operator} defines a
holomorphic map from $\Omega_{p,q}$ to
$\mathcal L(L^p(M,\nu),L^q(M,\nu))$. Its norm satisfies
\begin{equation}\label{eq:scale-operator-bound}
 \|\mathcal T_{n,R}(w,\omega)\|_{p\to q}
       \le c_*^{1/t}B_{t\sigma}^{1/t}e^{\sigma vn}.
\end{equation}
On every compact subset of $\Omega_{p,q}$, all fixed-order derivatives in $(w,\omega)$ have bounds $C_k e^{C_k n}$ uniform in $R$. At every fixed
$(w,\omega)\in\Omega_{p,q}$ and every $h\in L^p(M,\nu)$,
\begin{equation}\label{eq:scale-parameter-continuity}
 R\longmapsto\mathcal T_{n,R}(w,\omega)h
       \quad\hbox{is continuous in }L^q(M,\nu).
\end{equation}
For $\operatorname{Re}w<0$ and real $\omega$, this operator agrees
with $\mathcal R_R(e^w,\omega)^n$ wherever the latter is considered
as an $L^1$ operator. The resulting inserted record pairings extend
holomorphically across $w=0$ on this fixed neighborhood.
\end{theorem}
\begin{proof}
Change variables $y=(F_R^*)^n x$ in the $q$th power of the norm. The
absolute weight is at most $e^{\sigma N_{n,R}(x)}$, by
\eqref{eq:cumulative-record-domination}. H\"older's inequality with
$1/q=1/p+1/t$ and Theorem~\ref{thm:cumulative-return-tail} give
\[
 \|\mathcal T_{n,R}(w,\omega)h\|_q
 \le\left(\int_{Y_R^*}e^{t\sigma N_{n,R}}\dd\nu\right)^{1/t}
       \|h\|_p
 \le c_*^{1/t}B_{t\sigma}^{1/t}e^{\sigma vn}\|h\|_p.
\]
For $\sigma=0$ use $B_0=1$. This proves the operator bound, including
$p=\infty$.

On a compact subset of the stated domain choose a slightly larger
compact neighborhood for which $t\sigma'<c$. A derivative of order
$k$ in $(w,\omega)$ inserts a monomial in $N_{n,R}$ and the four
coordinates of $J_{n,R}$, bounded by $C_kN_{n,R}^k$. Every such
polynomial is absorbed by $e^{(\sigma'-\sigma)N_{n,R}}$. H\"older's
inequality and \eqref{eq:cumulative-exponential-moment} therefore bound
the derivative integrands and the Taylor remainders in operator norm.
The exponential power series on a sufficiently small polydisc converges
in that norm. This proves joint holomorphy, not just weak holomorphy of
individual pairings. The same estimates on the larger neighborhood,
or Cauchy's inequalities there, give the claimed derivative bounds.

We give the parameter-continuity argument because the underlying return
domains move. Fix $R_0$. Remove $\partial Y_{R_0}^*$, the forward and
backward collision singularities, and all finite collision iterates of
$\partial Y_{R_0}^*$. Each is a countable union of null curves, up to
singular endpoints. Backward recurrence holds almost everywhere. At a
remaining point of $Y_{R_0}^*$ the preceding $n$ returns use finitely many
nonsingular physical collisions; none is on a section boundary. The
implicit collision-root argument of
Theorem~\ref{thm:return-stability}, applied backwards, preserves this
finite itinerary for $R$ sufficiently near $R_0$. The backward starting
point and all flight lengths vary continuously; the lattice labels and
collision counts are fixed. At a remaining point outside $Y_{R_0}^*$
the section indicator stays zero. Thus
\eqref{eq:block-scale-operator} converges pointwise almost everywhere
when $h$ is bounded and continuous.

Choose $q'>q$ sufficiently close to $q$ so that $q'<p$ and
$t'\sigma<c$, where $t'=(1/q'-1/p)^{-1}$. The preceding bound in
$L^{q'}$ is uniform in nearby $R$. It gives uniform integrability of the
$q$th powers, so pointwise convergence implies convergence in $L^q$.
For $p<\infty$, approximate an arbitrary $h$ in $L^p$ by bounded smooth
functions and use the uniform $p\to q$ bound. For $p=\infty$, choose a
finite $p_0>q$ so large that $(1/q-1/p_0)^{-1}\sigma<c$. Approximation
in $L^{p_0}$ and its uniform $p_0\to q$ bound give the same conclusion
for $h\in L^\infty$. This avoids assuming norm density of smooth
functions in $L^\infty$.

In the damped region the two operators have the same pointwise formula,
by Theorem~\ref{thm:common-renewal}. For example, for $h\in L^p$ and
$b\in L^{q^\dagger}$, where $q^\dagger$ is the conjugate exponent of $q$,
\[
 c_*^{-1}\int b\,\mathcal T_{n,R}(w,\omega)h\dd\nu
 =\int_{Y_R^*}h(x)b((F_R^*)^nx)
       e^{wN_{n,R}(x)+i\omega\cdot J_{n,R}(x)}\dd\nu_R^*.
\]
This continuous pairing is holomorphic on $\Omega_{p,q}$. It agrees
with the damped renewal pairing there and hence is the asserted
extension across $w=0$.
\end{proof}

\subsection*{Compatibility of the realizations}
For real $\omega$ and $|z|<1$, the constructions fit the diagram
\[
 \begin{array}{ccc}
 \mathcal A_{R,\omega}&\longmapsto&\mathcal R_R(z,\omega)\\
 \Big\downarrow {\scriptstyle P_R(I-z\mathcal A)^{-1}P_R}
 &&\Big\downarrow {\scriptstyle (I-\mathcal R)^{-1}}\\
 \mathcal U_R(z,\omega)&=&\mathcal U_R(z,\omega).
 \end{array}
\]
The upper arrow is the first-return sum
\eqref{eq:schur-return-operator}; the bottom equality is
\eqref{eq:physical-renewal-identity}. On the common domain
$z=e^w$, $\operatorname{Re}w<0$, $\omega\in\mathbb R^4$, the second
compatibility is
\[
 \mathcal R_R(e^w,\omega)^n=\mathcal T_{n,R}(w,\omega),\qquad
 c_*^{-1}\langle b,\mathcal T_{n,R}h\rangle_\nu
 =\int h\,(b\circ(F_R^*)^n)e^{wN_{n,R}+i\omega\cdot J_{n,R}}
                        \dd\nu_R^*.
\]
The first identity is interpreted on $L^p\subset L^1$, with values in
$L^q$, when the scale operator is used; the second has the insertion
exponents of the theorem. Any anisotropic realization must reproduce
these pairings on its admissible test class. Such an additional
realization is not an arrow already established by this diagram.

\begin{remark}[Relation to the spectral and raw-density estimates]
The operator identities concern the genuine four-coordinate record,
and all their spaces live on the same collision probability space.
They specify a common realization with fixed exponential integrability
and strong parameter continuity. They do not assert operator-norm
continuity in $R$, quasi-compactness on $L^p$, or a holomorphic family of
endomorphisms of a single anisotropic space. At real frequencies the
induced $L^1$ operator is an isometry on its section, so its construction
alone gives no spectral gap. Likewise, maps $L^p\to L^q$ with $q<p$
cannot simply be iterated on one space outside the damped region; the
$n$-block operator above was defined directly.

The next analytical step is to establish an anisotropic realization
compatible with \eqref{eq:actual-record-pairing} and
\eqref{eq:physical-renewal-identity}, and to prove its phase
reconstruction and raw-residual bounds. The collision compensation
in Sections~\ref{sec:collision-covariance}--\ref{sec:functional-gaussian}
will establish the actual Gaussian covariance and functional limits
by a different route, without identifying an induced spectral family. The criteria in
Theorem~\ref{thm:phase-resolvent-criterion},
Proposition~\ref{prop:covariance-closure}, and
Lemma~\ref{lem:raw-residual-sum} continue to state these requirements
separately. In particular the common Lebesgue realization is not used
as a substitute for any of those estimates in the raw local limit.
\end{remark}
